\section{A claimed vanishing theorem for principal divisors}
\label{app:claimed-vanishing}

Mochizuki's Theorem~5 in version~8 of
\cite{Mochizuki2020} asserts that, for a commutative Noetherian local
ring $B$ and a nonzerodivisor $g$, the inclusion from finitely generated
projective $B/gB$-modules to finitely generated $B$-modules of projective
dimension at most one induces zero on $K$-theory. This statement reaches
the exact support in the present paper: Lemma~\ref{lem:local-ring} gives
$B=A$, $g=Y_0x-X_0y$, and $A/gA=\OO$, with $g$ a regular parameter.
Composing the claimed zero map with the inclusion into coherent
$A$-modules would annihilate $i_*\widetilde c$.

We give a counterexample to the functor assertion in
\cite[Lemma~13(III)]{Mochizuki2020}, which is used in that proof.
The counterexample already has both $B$ and $B/gB$ regular. It concerns
the supplied proof of the vanishing assertion; a failure of this
construction alone does not disprove the assertion itself. The proof
of Theorem~\ref{thm:main} above is independent of it.

\subsection*{The complexes and the proposed operation}

Let $B$ be a regular local ring and let $g$ be a nonzero regular
parameter. In particular $g$ is a nonzerodivisor in the maximal ideal,
and $1-g$ is a unit. The exact category $\mathcal C$ used in the cited
argument consists of two-term complexes of finite free $B$-modules,
in homological degrees one and zero, whose differential is injective
and whose cokernel is killed by $g$. Exact sequences are degreewise
exact. We use only the objects
\[
 (n,m)_B=
 \left[B^n\oplus B^m
 \xrightarrow{\operatorname{diag}(gI_n,I_m)}B^n\oplus B^m\right]
 \quad(n,m\geq0)
\]
and split exact sequences between them. Thus no classification of
other objects of $\mathcal C$ is needed for the example.

For $X=(1,1)_B$, put $D=\operatorname{diag}(g,1)$. Following the
block convention in \cite[Conventions~11, equations~(15)--(16)]{Mochizuki2020},
an encoded endomorphism
$M=\left(\begin{smallmatrix}a&b\\c&d\end{smallmatrix}\right)$ has
actual matrices
\begin{equation}\label{eq:mochizuki-chain-matrices}
 M_1=\begin{pmatrix}a&b\\gc&d\end{pmatrix},
 \qquad
 M_0=\begin{pmatrix}a&gb\\c&d\end{pmatrix}.
\end{equation}
These satisfy $DM_1=M_0D$. Conversely, the chain-map condition, together
with the fact that $g$ is a nonzerodivisor, gives this presentation.
Multiplying the actual matrices shows that composition of the encoded
matrices is
\begin{equation}\label{eq:mochizuki-matrix-composition}
 \begin{pmatrix}a'&b'\\c'&d'\end{pmatrix}
 \star
 \begin{pmatrix}a&b\\c&d\end{pmatrix}
 =
 \begin{pmatrix}
 a'a+gb'c&a'b+b'd\\
 c'a+d'c&gc'b+d'd
 \end{pmatrix}.
\end{equation}
Thus $M'\star M$ represents $M'\circ M$. Upper triangular means
$c=0$, and lower triangular means $b=0$; these terms refer to the
encoded matrices.

Write $\operatorname{Typ}(h)=[B\xrightarrow{h}B]$ for $h\in B$.
The assignments $\mu'_1,\mu'_2$ preceding Definition~12 in
\cite{Mochizuki2020} replace $(n,m)_B$ by
$\operatorname{Typ}(g)^{\oplus n}$ or
$\operatorname{Typ}(1)^{\oplus n}$, respectively. On a morphism they
retain its upper-left block in both chain degrees. On all morphisms
this is not a functor, because the upper-left entry of
\eqref{eq:mochizuki-matrix-composition} contains $gb'c$.
The cited argument therefore restricts the diagrams and their maps.

Here is the part of that restriction needed below. An $S_2$-object in
an exact category is an admissible short exact sequence
$X_{01}\rightarrowtail X_{02}\twoheadrightarrow X_{12}$, together
with the zero objects on its diagonal; morphisms are commuting maps
of these sequences. Definition~12 of \cite{Mochizuki2020} requires
upper triangular structure arrows and allows lower triangular
componentwise isomorphisms between the diagrams. Lemma~13(III) asserts
that $\mu'_1,\mu'_2$ induce functors on this modified $S$-construction.
The following example violates the naturality required of an image
morphism already in degree two.

\begin{lemma}\label{lem:triangular-naturality-failure}
For every regular local ring $B$ with a nonzero regular parameter $g$,
the upper-left-block assignments just described fail to send all
lower triangular isomorphisms between upper triangular $S_2$-objects
to morphisms of the target diagrams.
\end{lemma}

\begin{proof}
Take the encoded endomorphisms of $X=(1,1)_B$
\[
 U=\begin{pmatrix}1&1\\0&1\end{pmatrix},\qquad
 U'=\begin{pmatrix}1-g&1\\0&1\end{pmatrix},\qquad
 L=\begin{pmatrix}1&0\\1&1\end{pmatrix},\qquad
 L'=\begin{pmatrix}1&0\\1&1-g\end{pmatrix}.
\]
Their actual chain matrices are
\[
\begin{array}{c@{\qquad}c}
 U_1=\begin{pmatrix}1&1\\0&1\end{pmatrix}&
 U_0=\begin{pmatrix}1&g\\0&1\end{pmatrix}\\[6pt]
 U'_1=\begin{pmatrix}1-g&1\\0&1\end{pmatrix}&
 U'_0=\begin{pmatrix}1-g&g\\0&1\end{pmatrix}\\[6pt]
 L_1=\begin{pmatrix}1&0\\g&1\end{pmatrix}&
 L_0=\begin{pmatrix}1&0\\1&1\end{pmatrix}\\[6pt]
 L'_1=\begin{pmatrix}1&0\\g&1-g\end{pmatrix}&
 L'_0=\begin{pmatrix}1&0\\1&1-g\end{pmatrix}.
\end{array}
\]
In each degree the determinants of $U,U',L,L'$ are, respectively,
$1,1-g,1,1-g$. All are units, so these are chain isomorphisms.
The first two are upper triangular and the last two lower triangular.
A direct calculation with \eqref{eq:mochizuki-matrix-composition} gives
\[
 L'\star U=\begin{pmatrix}1&1\\1&1\end{pmatrix}=U'\star L.
\]
Equivalently, both degree-one composites are
$\left(\begin{smallmatrix}1&1\\g&1\end{smallmatrix}\right)$ and both
degree-zero composites are
$\left(\begin{smallmatrix}1&g\\1&1\end{smallmatrix}\right)$.
Hence the diagram
\[
\begin{tikzcd}[column sep=large]
 X\arrow[r,"U"]\arrow[d,"L"'] &
 X\arrow[r]\arrow[d,"L'"] &0\arrow[d,equal]\\
 X\arrow[r,"U'"'] &X\arrow[r]&0
\end{tikzcd}
\]
commutes. Each row is split exact since its first arrow is an
isomorphism. Besides that first arrow, all structure arrows of the
corresponding $S_2$-object are zero or identity maps and are therefore
upper triangular. The vertical maps $(L,L',\operatorname{id}_0)$
are lower triangular isomorphisms. Thus the diagram is a morphism in
the proposed source category.

Both upper-left entries of $L,L'$ are $1$, whereas those of $U,U'$ are
$1,1-g$. Either assignment would therefore send the left square to
\[
\begin{tikzcd}[column sep=large,row sep=large]
 \operatorname{Typ}(h)\arrow[r,"1"]\arrow[d,"1"']&
 \operatorname{Typ}(h)\arrow[d,"1"]\\
 \operatorname{Typ}(h)\arrow[r,"1-g"']&\operatorname{Typ}(h),
\end{tikzcd}
\qquad h=g\text{ or }h=1.
\]
Its two composites are $1$ and $1-g$ in both degrees. They are unequal
because $g\ne0$. The proposed image is therefore not a morphism of
$S_2$-diagrams.
\end{proof}

The calculation identifies the precise missing condition: preservation
of upper--upper and lower--lower compositions does not imply preservation
of mixed upper--lower commuting squares. In version~8 the encoded
composition and assignments appear on page~9, Lemma~13(III) on page~10,
and the subsequent additivity argument on page~11. That argument uses
the asserted functors in Lemma~13(IV) and in the proof of assertion
$(\beta)$. Although the two projected composites can be chain homotopic,
the target category requires strictly commuting natural transformations.
The later homotopy comparison does not define replacement functors on
this square. A construction that supplies the required higher
compatibilities and additivity would be an additional argument.
